\section{References}

\begingroup
\raggedright

Ancient and patristic citations in the text follow conventional titles and loci
keyed to the public-domain editions listed below. Witnesses consulted through
secondary summaries rather than standalone full editions are marked as reported.

\subsection*{Primary sources: Gregory the Great}

\begin{itemize}
\item Gregory the Great, \work{The Book of Pastoral Rule} (\work{Liber Regulae
Pastoralis}) and \work{Selected Epistles} (Books I--VIII), trans.\ James
Barmby, with Prolegomena and notes, in \textit{A Select Library of Nicene and
Post-Nicene Fathers of the Christian Church}, Series 2, vol.~12, ed.\ Philip
Schaff and Henry Wace (New York: Christian Literature Company, 1895; CCEL
unrevised printing). Epistle numbers follow the Benedictine/Maurist
arrangement printed by Barmby. Public domain.
\item Gregory the Great, \work{Morals on the Book of Job} (\work{Moralia in
Iob}), trans.\ James Bliss and Charles Marriott, Library of the Fathers
(Oxford: John Henry Parker, 1844--1850). Inspected at the Preface, Epistle
Dedicatory to Leander of Seville, Book I opening, and Book XIV.56. Public
domain.
\item Gregory the Great, \work{Homiliae XL in Evangelia} and \work{Oratio ad
plebem de mortalitate}, Latin text from the Maurist/Benedictine edition in
J.-P.\ Migne, \textit{Patrologia Latina}, vol.~76 (Paris: Garnier, 1878),
cols.\ 1075--1314 (including the Preface to Secundinus, Homilies 1.1, 1.5, 8.1,
20, 37.7--9, 38.15--16, 40.7, 40.11--12, and \work{Oratio ad plebem de
mortalitate}, cols.\ 1311--1314). Public domain.
\item Gregory the Great, \work{Homiliae in Hiezechihelem Prophetam}, Latin text
from J.-P.\ Migne, \textit{Patrologia Latina}, vol.~76 (Paris, 1849), inspected
at Book II, Homilies 6, 7, and 10. Public domain.
\item Gregory the Great, \work{Dialogi} (Books I--IV) and \work{Registrum
Epistolarum} Books IX--XIV. Cited at reported level through James Barmby's
Prolegomena and notes in NPNF2-12, Adrian Fortescue (1922), and editorial
cross-references in Migne PL 76.
\end{itemize}

\subsection*{Contemporary external and monastic witnesses}

\begin{itemize}
\item Gregory of Tours, \work{The History of the Franks} (\work{Historia
Francorum}), Book X.1, trans.\ O.\ M.\ Dalton, vol.~2 (Oxford: Clarendon
Press, 1927), pp.~425--428, collated with the Latin critical edition of
W.\ Arndt, \textit{Monumenta Germaniae Historica, Scriptores rerum
Merovingicarum}, vol.~1 (Hanover, 1884), pp.~406--409. Public domain.
\item Benedict of Nursia, \work{Regula Benedicti}, Latin text consulted for
sixth-century Italian monastic context beside Gregory's \work{Dialogi} Book II.
Public domain.
\end{itemize}

\subsection*{Liturgical and official ecclesial witnesses}

\begin{itemize}
\item H.\ A.\ Wilson, ed., \work{The Gregorian Sacramentary under Charles the
Great, Edited from Three MSS.\ of the Ninth Century} (\work{Hadrianum}), Henry
Bradshaw Society, vol.~49 (London: Harrison and Sons, 1915). Inspected at
Wilson's Introduction, the Latin \latin{Hanc igitur} of the Roman Canon, the
Mass \latin{Natale Sancti Gregorii Papae} (p.~25), and the Carolingian
\latin{Hucusque} preface (pp.~145--146). Public domain.
\item J.-P.\ Migne, ed., \textit{Patrologia Latina}, vol.~78 (\work{Sancti
Gregorii Magni Opera Omnia}, vol.~4, Maurist/Benedictine edition, Paris).
Inspected at the prefaces to the \work{Liber Sacramentorum} (including Pope
Hadrian I's letter to Charlemagne, cols.\ 23--26) and the \work{Ad
Antiphonarium S. Gregorii Praefatio} with Pierre de Goussainville's preface
quoting John the Deacon, \work{Vita S. Gregorii} II.6--9 (cols.\ 649--651).
Public domain.
\item \work{Martyrologium Romanum}, editio novissima under Leo XIII (Rome: Ex
typographia polyglotta S.\ C.\ de Propaganda Fide, 1902), inspected at 12 March
(p.~36) and 3 September (p.~132). Public domain.
\end{itemize}

\subsection*{Early lives, later tradition, and modern scholarship}

\begin{itemize}
\item Bede the Venerable, \work{Historia Ecclesiastica Gentis Anglorum}
(I.23--32, II.1, V.24); anonymous Whitby \work{Vita Gregorii} (early 8th c.);
Paul the Deacon, \work{Vita Gregorii} and \work{Historia Langobardorum}
(III.24, 33, 36; IV.8, 29); John the Deacon, \work{Vita S. Gregorii} (Books
I--IV); and Pope Boniface VIII's 1298 decretal on the four Latin Doctors. Cited
at reported level through James Barmby's Prolegomena and notes (1895),
H.\ A.\ Wilson (1915), Adrian Fortescue (1922), and Migne PL 76 and PL 78.
\item Adrian Fortescue, \work{The Mass: A Study of the Roman Liturgy} (London:
Longmans, Green, 1922), pp.~172 and 234--235, documenting Gregory's liturgical
reforms (\latin{Kyrie eleison}, \latin{Hanc igitur}, and \latin{Pater noster}),
quoting \work{Ep.}\ IX.12 on Roman \latin{Kyrie} and \latin{Christe eleison}
usage, and citing the Council of Vaison (529) against John the Deacon. Public
domain.
\end{itemize}

\endgroup
